\documentclass[10pt,a4paper]{amsart}
\usepackage[margin=19mm]{geometry}
\usepackage{amsmath,amssymb,mathtools}
\usepackage[scr=rsfs,cal=euler]{mathalfa}
\usepackage{xfrac}
\usepackage[unicode,hidelinks,bookmarks=false]{hyperref}
\newcommand{\ie}{i.e.\ }
\newcommand{\cf}{cf.\ }
\newcommand{\resp}{respectively\ }
\newcommand{\Subsection}{\subsection}
\providecommand{\nobreakdash}{\nobreak\hbox{-}}
\setlength{\emergencystretch}{2em}
\multlinegap0pt
\usepackage{latexsym,amscd}
\newtheorem{theorem}{Theorem}[section]
\newtheorem{example}{Example}[section]
\usepackage{mathrsfs}
\usepackage{mathtools}
\usepackage{youngtab}
\usepackage{enumitem}
\newcommand{\Int}{\operatorname{Int}}
\newcommand{\Orb}{\operatorname{Orb}}
\newcommand{\Deck}{\operatorname{Deck}}
\newcommand{\Ext}{\operatorname{Ext}}
\newcommand{\Stab}{\operatorname{Stab}}
\newcommand{\Aut}{\operatorname{Aut}}
\newcommand{\aut}{\operatorname{aut}}
\newcommand{\Tr}{\operatorname{Tr}}
\newcommand{\Del}{\operatorname{Del}}
\newcommand{\ord}{\operatorname{ord}}
\newcommand{\End}{\operatorname{End}}
\newcommand{\Sym}{\operatorname{Sym}}
\begin{document}
\title[The Graph Algebra I: Representation-Theoretic Structure]{The Graph Algebra I: Representation-Theoretic Structure}
\author{Leonid Bedratyuk}
\date{}
\maketitle
\noindent\emph{Source and licence.} The Graph Algebra I: Representation-Theoretic Structure. Version arXiv:2606.29558v1 (CC BY 4.0). This material is distributed under the Creative Commons Attribution 4.0 International licence, http://creativecommons.org/licenses/by/4.0/, as recorded in the arXiv OAI licence metadata for this exact version and shown on the abstract page https://arxiv.org/abs/2606.29558v1. Full rights notice: licenses/arxiv-2606.29558-CC-BY-4.0.txt.\par\smallskip
\noindent\textit{Editorial scope.} Counted unit: the original \texttt{theorem} environment \texttt{-} at source lines 1229--1250 together with its complete proof at lines 1252--1274; the delivered document keeps source lines 606--1392 so that every referenced label and every citation resolves inside it. Native editable source is the paper's own concatenated TeX; the AMS wrapper is editorial formatting only. No publisher class, upstream script or unrelated artwork is redistributed.
\begin{theorem}\label{complete_reducibility_sl2}
  Every finite-dimensional $\mathfrak{sl}_2$-module over a field of
  characteristic zero is isomorphic to a direct sum of irreducible
  submodules:
  $$
    V \cong \bigoplus_{d \ge 0} m_d\, V_d,
  $$
  where the multiplicities $m_d \in \mathbb{Z}_{\ge 0}$, almost all of which
  are zero, are uniquely determined.
\end{theorem}

\medskip
The decomposition of a module~$V$ into irreducible submodules is practically
found using elements of the kernel $\ker e_-$ of the operator $e_-$.
Each nonzero vector $v_0 \in \ker e_- \cap V_{d}$ generates an irreducible
submodule
$\langle v_0, e_+v_0, e_+^2v_0, \dots \rangle \cong V_d$;
the multiplicity $m_d$ is equal to the dimension of the corresponding weight
subspace:
$$
  m_d = \dim(\ker e_- \cap V_{d}).
$$
Thus, the complete information about the decomposition is encoded in the
structure of $\ker e_-$.

\medskip

In the next section we define on the graph algebra~$\mathcal{A}$ the
structure of an $\mathfrak{sl}_2$-module.








%%%%%%%%%%%%%%%%%%%%%%%%%%%%%%%
%%%%%%%%%%%%%%%%%%%%%%
\section{The graph algebra $\mathcal A$ as an $\mathfrak{sl}_2$-module}
%%%%%%%%%%%%%%%%%%%%%%%%%%%%%%%%%%%%%
%%%%%%%%%%%%%%%%%%%%%%%%%%%%%%%%%%%%

In this section we show that the graph algebra $\mathcal A$ is naturally
endowed with the structure of a finite-dimensional $\mathfrak{sl}_2$-module.
The idea of this construction is quite simple: if a squarefree monomial
$\boldsymbol{x}_S$ is interpreted as a graph with edge set $S\subseteq[m]$,
then the natural operations on such monomials are adding one new edge,
deleting one existing edge, and recording the number of edges.
These three operations lead to operators which, as it turns out, satisfy the
standard commutation relations of the algebra $\mathfrak{sl}_2$.
As a result, the rank structure of the algebra $\mathcal A$ is compatible
with its weight decomposition, and the problem of describing $\mathcal A$
is reduced to classical questions in the representation theory of
$\mathfrak{sl}_2$.

\subsection{The action of the algebra $\mathfrak{sl}_2$ on $\mathcal {A}$}

Consider on $\mathcal {A}$ three linear operators defined on basis monomials as follows:
\begin{gather*}
D_+(\boldsymbol{x}_S)=\sum_{e\in\overline S}\boldsymbol{x}_{S\cup\{e\}}, D_+(\boldsymbol{x}_{[m]})=0,\\
D_0(\boldsymbol{x}_S)=-(m-2|S|)\,\boldsymbol{x}_S,\\
D_-(\boldsymbol{x}_S)=\sum_{e\in S}\boldsymbol{x}_{S\setminus\{e\}},D_-(\boldsymbol{x}_\varnothing)=0.
\end{gather*}
%The operator $D_+(\boldsymbol{x}_S)$ is called the \emph{raising operator} and is the sum of all monomials obtained from $\boldsymbol{x}_S$ by adding one new edge.
%The operator $D_-(\boldsymbol{x}_S)$ is called the \emph{lowering operator} and is the sum of all monomials obtained from $\boldsymbol{x}_S$ by deleting one edge.
%The operator $D_0$ is called the \emph{weight operator}; it acts diagonally in the basis $\{\boldsymbol{x}_S\}$ and records the number of edges in the graph $S$.

All three operators are compatible with the rank structure on $\mathcal A:$
$$
D_+:A_k\to A_{k+1},\qquad
D_-:A_k\to A_{k-1},\qquad
D_0:A_k\to A_k.
$$
%The operators $D_+$ and $D_-$ are locally nilpotent operators on $\mathcal A.$.

We now show that these three operators $D_+,D_0,D_-$ define the standard
action of the algebra $\mathfrak{sl}_2$ on $\mathcal A$.
\begin{theorem}
The operators $D_+,D_0,D_-$ satisfy the relations
\begin{equation}\label{sl2_on_A_relations}
[D_+,D_-]=D_0,\qquad
[D_0,D_+]=2D_+,\qquad
[D_0,D_-]=-2D_-,
\end{equation}
and define on the algebra $\mathcal A$ the structure of an
$\mathfrak{sl}_2$-module.
\end{theorem}

\begin{proof}
It is enough to check the relations \eqref{sl2_on_A_relations} on basis
monomials $\boldsymbol{x}_S, |S|=k$, and then use linearity.

\smallskip

First we verify the relation
$
[D_0,D_+]=2D_+.
$
Every monomial in the sum
$$
D_+(\boldsymbol{x}_S)=\sum_{e\in\overline S}\boldsymbol{x}_{S\cup\{e\}},
$$
has degree $k+1$, and therefore
$$
D_0(D_+(\boldsymbol{x}_S))=-(m-2(k+1))D_+(\boldsymbol{x}_S).
$$
On the other hand,
$$
D_+(D_0(\boldsymbol{x}_S))=D_+\bigl(-(m-2k)\boldsymbol{x}_S\bigr)=-(m-2k)D_+(\boldsymbol{x}_S).
$$
Hence
$$
[D_0,D_+](\boldsymbol{x}_S)
=
D_0(D_+(\boldsymbol{x}_S))-D_+(D_0(\boldsymbol{x}_S))
=
\bigl(-(m-2k-2)+(m-2k)\bigr)D_+(\boldsymbol{x}_S)
=
2D_+(\boldsymbol{x}_S).
$$

The relation $[D_0,D_-]=-2D_-$ is verified similarly:
$$
[D_0,D_-](\boldsymbol{x}_S)
=
D_0(D_-(\boldsymbol{x}_S))-D_-(D_0(\boldsymbol{x}_S))
=
\bigl(-(m-2k+2)+(m-2k)\bigr)D_-(\boldsymbol{x}_S)
=
-2D_-(\boldsymbol{x}_S).
$$

It remains to verify the last relation
$
[D_+,D_-]=D_0.
$
We have
$$
D_-D_+(\boldsymbol{x}_S)
=
\sum_{f\in\overline S}D_-(\boldsymbol{x}_{S\cup\{f\}}).
$$
For a fixed $f\in\overline S$ we obtain
$$
D_-(\boldsymbol{x}_{S\cup\{f\}})
=
\sum_{e\in S\cup\{f\}}\boldsymbol{x}_{(S\cup\{f\})\setminus\{e\}}
=
\boldsymbol{x}_S+\sum_{e\in S}\boldsymbol{x}_{(S\setminus\{e\})\cup\{f\}}.
$$
Therefore,
$$
D_-D_+(\boldsymbol{x}_S)
=
(m-k)\boldsymbol{x}_S+\sum_{f\in\overline S}\sum_{e\in S}\boldsymbol{x}_{(S\setminus\{e\})\cup\{f\}}.
$$

On the other hand,
$$
D_-(\boldsymbol{x}_S)=\sum_{e\in S}\boldsymbol{x}_{S\setminus\{e\}},
$$
and hence
$$
D_+D_-(\boldsymbol{x}_S)
=
\sum_{e\in S}D_+(\boldsymbol{x}_{S\setminus\{e\}}).
$$
For a fixed $e\in S$, the set of edges absent from $S\setminus\{e\}$ is
$$
[m]\setminus (S\setminus\{e\})=\overline S\cup\{e\},
$$
and therefore
\begin{equation}\label{r2}
D_+(\boldsymbol{x}_{S\setminus\{e\}})
=
\boldsymbol{x}_S+\sum_{f\in\overline S}\boldsymbol{x}_{(S\setminus\{e\})\cup\{f\}}.
\end{equation}
It follows that
$$
D_+D_-(\boldsymbol{x}_S)
=
k\boldsymbol{x}_S+\sum_{e\in S}\sum_{f\in\overline S}\boldsymbol{x}_{(S\setminus\{e\})\cup\{f\}}.
$$

The double sums in the last two formulas coincide, since they are taken over
the same pairs
$$
(e,f)\in S\times\overline S.
$$
Thus, after subtraction, they cancel each other, and we obtain
$$
[D_+,D_-](\boldsymbol{x}_S)
=
D_+D_-(\boldsymbol{x}_S)-D_-D_+(\boldsymbol{x}_S)
=
(2k-m)\boldsymbol{x}_S
=
-(m-2 |S|)\boldsymbol{x}_S
=
D_0(\boldsymbol{x}_S).
$$

Hence all relations \eqref{sl2_on_A_relations} hold on the basis, and
therefore on the whole algebra $\mathcal A$.
\end{proof}
The combinatorial meaning of the operators $D_+,D_-,D_0$ is completely
transparent: $D_+$ runs through all one-edge supergraphs of a given graph,
$D_-$ runs through all its one-edge subgraphs, and $D_0$ records the number
of edges.
It is precisely this interpretation that makes the constructed
$\mathfrak{sl}_2$-structure natural from the point of view of graph theory.

\subsection{Primitive elements}

The elements of the algebra $\mathcal A$ which are annihilated by the operator
$D_-$ form the subspace
$$
\ker D_-=\{f\in\mathcal A:\ D_-f=0\}.
$$
The kernel of the operator $D_-$ is a rank subspace in $\mathcal A$
$$
\ker D_-= P_0+P_1+ \cdots + P_m,
$$
where
$$
P_k=\ker\left(D_-:A_k\longrightarrow A_{k-1}\right),
\qquad A_{-1}=0.
$$
Elements of the space $P_k$ will be called \emph{primitive elements of degree}
$k$.

The following theorem shows that primitive elements are initial vectors of
irreducible $\mathfrak{sl}_2$-chains and that such elements can appear only
in the lower half of the algebra.
\begin{theorem}\label{primitive_chain}
Let $0\neq z\in P_k$. If $k\le \lfloor m/2\rfloor$, then the elements
$$
z,\ D_+(z),\ D_+^2(z),\dots,\ D_+^{m-2k}(z),
$$
generate an irreducible $\mathfrak{sl}_2$-module isomorphic to $V_{m-2k}$.
In particular,
$$
D_+^{m-2k}(z)\neq 0,
\qquad
D_+^{m-2k+1}(z)=0.
$$
If $k>\lfloor m/2\rfloor$, then $P_k=0$.
\end{theorem}

\begin{proof}
Since $\mathcal A$ is a finite-dimensional $\mathfrak{sl}_2$-module, by the
complete reducibility theorem it is a direct sum of irreducible
$\mathfrak{sl}_2$-modules. Let $0\neq z\in P_k$. Then $z\in A_k$,
$D_-z=0$, and the weight of $z$ is
$$
-(m-2k).
$$
Thus, $z$ is a lowest weight vector in the submodule that it generates.
Therefore this submodule is isomorphic to the standard irreducible module
$V_{m-2k}$. It follows that the chain
$$
z,\ D_+z,\ D_+^2z,\dots,\ D_+^{m-2k}z,
$$
has length $m-2k+1$, and hence
$$
D_+^{m-2k}(z)\neq 0,
\qquad
D_+^{m-2k+1}(z)=0.
$$

If $k>m/2$, then the weight $2k-m$ is positive. But a nonzero vector
annihilated by the operator $D_-$ in a finite-dimensional
$\mathfrak{sl}_2$-module can only be a lowest weight vector of some
irreducible component, and the lowest weights of irreducible modules are of
the form $-d\le 0$. Hence, for $k>m/2$, there are no such vectors, that is,
$P_k=0$.
\end{proof}

The following theorem gives the dimension of the vector space of primitive
elements $P_k$.

\begin{theorem}\label{primitive_dimensions}
For each $k=0,1,\dots,m$ we have
$$
\dim P_k=
\begin{cases}
\displaystyle \binom{m}{k}-\binom{m}{k-1},
& 0\le k\le \left\lfloor \dfrac m2\right\rfloor,\\[8pt]
0,
& \left\lfloor \dfrac m2\right\rfloor<k\le m.
\end{cases}
$$
\end{theorem}

\begin{proof}
Since $\mathcal A$ is a finite-dimensional $\mathfrak{sl}_2$-module over a
field of characteristic zero, it decomposes into a direct sum of irreducible
$\mathfrak{sl}_2$-modules.

In every irreducible module generated by a primitive element of degree $j$,
there is exactly one vector in each degree
$$
j,\ j+1,\dots,\ m-j,
$$
and there are no vectors in other degrees. Moreover, each such chain has a
unique initial vector annihilated by the operator $D_-$, while the chains that
started in smaller degrees have already been counted in the preceding
subspaces. Hence, if $k\le \lfloor m/2\rfloor$, the difference
$$
\dim A_k-\dim A_{k-1},
$$
counts exactly the number of new irreducible chains that start in degree
$k$. This number is equal to $\dim P_k$, because each such chain has exactly
one primitive vector in its initial degree.
Therefore
$$
\dim P_k
=
\dim A_k-\dim A_{k-1}
=
\binom{m}{k}-\binom{m}{k-1}.
$$
For $k>\lfloor m/2\rfloor$, the equality $P_k=0$ has already been established
in Theorem~\ref{primitive_chain}.
\end{proof}

As a consequence, for $1\le k\le \lfloor m/2\rfloor$ the operator
$
D_-
$
is surjective. Indeed,
$$
\operatorname{rank} D_-|_{A_k}
=
\dim A_k-\dim P_k
=
\binom{m}{k}
-
\left(\binom{m}{k}-\binom{m}{k-1}\right)
=
\binom{m}{k-1}
=
\dim A_{k-1}.
$$

Thus, for $k\le \lfloor m/2\rfloor$, the elements of the space $P_k$ are
precisely those elements in the space $A_k$ which are not obtained by the
action of the operator $D_+$ on elements of $A_{k-1}$.
It follows from the irreducible $\mathfrak{sl}_2$-decomposition that the
operator $D_+$ is injective on subspaces below the middle level. Therefore,
for $k\le \lfloor m/2\rfloor$, we have the direct sum
$$
A_k=P_k\oplus D_+(A_{k-1}).
$$

Iterating this decomposition, we obtain a decomposition of $A_k$ through
primitive spaces of smaller degrees:
\begin{equation} \label{P_to_A}
A_k
=
\bigoplus_{j=0}^{\min(k,m-k)}
D_+^{\,k-j}P_j.
\end{equation}

In other words, every element of $A_k$ is a sum of contributions from those
$\mathfrak{sl}_2$-chains which start in some primitive space $P_j$ and reach
$A_k$. The condition
$j\le \min(k,m-k)$ precisely means that the chain generated by $P_j$ contains
a component of degree $k$.

\medskip

Note that the operators $D_-$ and $D_+$ are not derivations of the graph
algebra. For the operator $D_-$ this is seen from the following formula:
$$
D_-(\boldsymbol{x}_S\boldsymbol{x}_T)
=
D_-(\boldsymbol{x}_S)\boldsymbol{x}_T
+
\boldsymbol{x}_S D_-(\boldsymbol{x}_T)
-
2|S\cap T|\,\boldsymbol{x}_{S\cup T}
+
\sum_{e\in S\cap T}
\boldsymbol{x}_{(S\cup T)\setminus\{e\}}.
$$
However, if $S\cap T=\varnothing$, then
\begin{equation}\label{Leibniz}
D_-(\boldsymbol{x}_S\boldsymbol{x}_T)=D_-(\boldsymbol{x}_S)\boldsymbol{x}_T+\boldsymbol{x}_SD_-(\boldsymbol{x}_T).
\end{equation}
Thus, the operator $D_-$ behaves as a derivation only on products of monomials
with disjoint supports.


\medskip

Let us describe all primitive elements for $n=4$.

\begin{example}
{\rm
By Theorem~\ref{primitive_dimensions} we have
$$
\dim P_0=1,\qquad
\dim P_1=5,\qquad
\dim P_2=9,\qquad
\dim P_3=5,
$$
and
$$
P_k=0,\qquad k>3.
$$
Let us describe these vector spaces of primitive elements explicitly.

\medskip

\noindent\textbf{Degree $0$.}
Since $A_0=\langle 1\rangle$ and $D_-(1)=0$, we have
$
P_0=\langle 1\rangle.
$
The element $1$ has weight $-6$ and generates an irreducible
$\mathfrak{sl}_2$-submodule isomorphic to $V_6$.

\medskip

\noindent\textbf{Degree $1$.}
On basis monomials of degree $1$ we have
$
D_-(x_i)=1,\qquad i=1,\dots,6,
$
and therefore
$$
P_1=\left\{\sum_{i=1}^6 a_i x_i:\ \sum_{i=1}^6 a_i=0\right\}.
$$
A convenient basis of this space is
$$
P_1=
\left\langle
x_1-x_2,\ x_1-x_3,\ x_1-x_4,\ x_1-x_5,\ x_1-x_6
\right\rangle.
$$
Every nonzero element of $P_1$ generates an irreducible module of type $V_4$.

\medskip

\noindent\textbf{Degree $2$.}
The space $P_2=\ker(D_-:A_2\to A_1)$ has dimension
$$
\dim P_2=\binom62-\binom61=9.
$$
Its basis can be chosen in the form of products of linear differences:
$$
\begin{aligned}
f_1&=(x_1-x_2)(x_3-x_4),&
f_2&=(x_1-x_2)(x_3-x_5),&
f_3&=(x_1-x_2)(x_3-x_6),\\
f_4&=(x_1-x_3)(x_2-x_4),&
f_5&=(x_1-x_3)(x_2-x_5),&
f_6&=(x_1-x_3)(x_2-x_6),\\
f_7&=(x_1-x_4)(x_2-x_5),&
f_8&=(x_1-x_4)(x_2-x_6),&
f_9&=(x_1-x_5)(x_2-x_6).
\end{aligned}
$$
Their primitiveness follows immediately from \eqref{Leibniz}, and their
linear independence follows from the fact that each $f_i$ contains a monomial
which does not occur in the expansion of any other element in the list.
Every nonzero element of $P_2$ generates an irreducible
$\mathfrak{sl}_2$-module of type $V_2$.

\medskip

\noindent\textbf{Degree $3$.}
Primitive elements of degree $3$ have weight $0$ and generate trivial
$\mathfrak{sl}_2$-modules of type $V_0$. The dimension of this space is
$$
\dim P_3=\binom63-\binom62=5.
$$
A convenient basis is given by five products of three linear differences:
$$
\begin{aligned}
g_1&=(x_1-x_2)(x_3-x_4)(x_5-x_6),\\
g_2&=(x_1-x_2)(x_3-x_5)(x_4-x_6),\\
g_3&=(x_1-x_3)(x_2-x_4)(x_5-x_6),\\
g_4&=(x_1-x_3)(x_2-x_5)(x_4-x_6),\\
g_5&=(x_1-x_4)(x_2-x_5)(x_3-x_6).
\end{aligned}
$$
Since this is a trivial $\mathfrak{sl}_2$-module, for every $g\in P_3$ we
also have
$
D_+(g)=0.
$

Thus, for $n=4$ the primitive spaces have the form
$$
P_0\cong \mathbb C,\qquad
\dim P_1=5,\qquad
\dim P_2=9,\qquad
\dim P_3=5,
$$
and the corresponding decomposition of the algebra as an
$\mathfrak{sl}_2$-module is
$$
\mathcal A
\cong
V_6\oplus 5V_4\oplus 9V_2\oplus 5V_0.
$$
Checking the dimension gives
$$
\dim\mathcal A
=
1\cdot 7+5\cdot 5+9\cdot 3+5\cdot 1
=
2^6.
$$
}
$\triangle$ \end{example}

%In the next subsection we will obtain the decomposition of $\mathcal A$ into irreducible $\mathfrak{sl}_2$-modules for arbitrary $n.$

In the general case, the following statement holds:

\begin{theorem}\label{irred_decomp_A}
The algebra $\mathcal A$, as an $\mathfrak{sl}_2$-module, decomposes into a
direct sum of irreducible modules
\begin{equation}
\mathcal A
\cong
\bigoplus_{k=0}^{\lfloor m/2\rfloor}
\mu_k\,V_{m-2k},
\end{equation}
where
\begin{equation}
\mu_k
=
\dim P_k
=
\binom{m}{k}-\binom{m}{k-1},
\qquad
0\le k\le \left\lfloor\frac m2\right\rfloor, \quad \binom{m}{-1}=0.
\end{equation}
\end{theorem}

The proof follows immediately from Theorem~\ref{complete_reducibility_sl2}
and Theorem~\ref{primitive_dimensions}.

Thus, we have obtained the complete decomposition of the graph algebra
$\mathcal A$ into irreducible $\mathfrak{sl}_2$-modules.
As a by-product, passing to dimensions, we obtain the  binomial
identity
$$
\sum_{k=0}^{\lfloor m/2\rfloor}
\left(\binom{m}{k}-\binom{m}{k-1}\right)(m-2k+1)=2^m.
$$


%%%%%%%%%%%%%%%%%%%%%%%%%%%
\subsection{Involution}
%%%%%%%%%%%%%%%%%%%%%%%%%%%%%%%%%%%

Define the linear operator
\begin{equation}\label{involution}
\omega:\mathcal A\to\mathcal A,
\qquad
\omega(\boldsymbol{x}_S)=\boldsymbol{x}_{\overline S}.
\end{equation}
It is easy to check that the operator $\omega$ is a linear involution:
$
\omega^2=\mathrm{Id}_{\mathcal A},
$
and moreover
$
\omega(A_k)=A_{m-k}.
$

The following statement holds.
\begin{theorem}
The operators $D_+$ and $D_-$ are related by
$$
D_+=\omega D_-\omega.
$$
\end{theorem}

\begin{proof}
It is enough to check the statement on basis monomials $\boldsymbol{x}_S$.
We have
$$
D_-(\omega(\boldsymbol{x}_S))=D_-(\boldsymbol{x}_{\overline S})
=
\sum_{e\in \overline S}\boldsymbol{x}_{\overline S\setminus\{e\}}.
$$
Now apply $\omega$ to each summand:
$$
\omega\!\left(\boldsymbol{x}_{\overline S\setminus\{e\}}\right)
=
\boldsymbol{x}_{[m]\setminus(\overline S\setminus\{e\})}
=
\boldsymbol{x}_{S\cup\{e\}}.
$$
Therefore
$$
\omega D_-\omega(\boldsymbol{x}_S)
=
\sum_{e\in \overline S}\boldsymbol{x}_{S\cup\{e\}}
=
D_+(\boldsymbol{x}_S).
$$
Since both operators are linear, we obtain
$
D_+=\omega D_-\omega
$
on the whole of $\mathcal A$.
\end{proof}

Let
$$
\mathcal A^{\mathfrak{sl}_2}=\ker D_-\cap \ker D_+,
$$
be the vector space of $\mathfrak{sl}_2$-invariants, that is, the elements of
$\mathcal A$ which generate trivial $\mathfrak{sl}_2$-modules.


\begin{theorem}
The following statements hold.
\begin{enumerate}
\item[\textnormal{(i)}] The graph-complement involution preserves the space of
$\mathfrak{sl}_2$-invariants:
$$
\omega\bigl(\mathcal A^{\mathfrak{sl}_2}\bigr)
=
\mathcal A^{\mathfrak{sl}_2}.
$$

\item[\textnormal{(ii)}]
The following equality holds:
$$
\mathcal A^{\mathfrak{sl}_2}=
\begin{cases}
P_{m/2}, \quad \text{if $m$ is even},\\
0, \quad \text{if $m$ is odd}.
\end{cases}
$$
\end{enumerate}
\end{theorem}

\begin{proof}
If $f\in\mathcal A^{\mathfrak{sl}_2}$, then by the properties of the
involution we have
$$
D_+(\omega f)=\omega D_-f=0,
\qquad
D_-(\omega f)=\omega D_+f=0.
$$
Hence
$$
\omega f\in\mathcal A^{\mathfrak{sl}_2},
$$
which proves part \textnormal{(i)}.
To prove part \textnormal{(ii)}, note that $\mathfrak{sl}_2$-invariants have
weight zero. This is possible only under the condition
$
m-2j=0.
$
If $m$ is odd, such a $j$ does not exist, and therefore
$\mathcal A^\mathfrak{sl}_2=0$.
If $m$ is even, then all invariants lie in $A_{m/2}$, and therefore
$\mathcal A^{\mathfrak{sl}_2}=P_{m/2}$.
\end{proof}


\medskip

It is worth noting that the constructed $\mathfrak{sl}_2$-structure is a
graph-theoretic realization of the standard mechanism of raising and lowering
operators on Boolean and related graded posets. Unlike general rank-theoretic
applications of this mechanism, in our case the operators $D_+$ and $D_-$ have
a direct graph-theoretic interpretation as adding and deleting one edge.
In the next section we show that this structure is compatible with the action
of the pair group $\Gamma$ and therefore passes to the space of graph
invariants.



%%%%%%%%%%%%%%%%%%%%%%%%%%%%%%%%%%%%%%%%%

\section{The graph algebra $\mathcal A$ as a 
$\Gamma$-module}

%%%%%%%%%%%%%%%%%%%%%%%%%%%%%%%%%%%%%%%%%

In the previous section, the graph algebra $\mathcal A$ was considered as an
$\mathfrak{sl}_2$-module. We now consider the structure on $\mathcal A$
defined by the action of the pair group
$
\Gamma.
$
Unlike the $\mathfrak{sl}_2$-structure, where the decomposition of the algebra
$\mathcal A$ has a simple explicit form, the decomposition with respect to the
group
$
\Gamma
$
is much more complicated. In order to describe it, we first consider the
action of the larger symmetric group
$
S_m
$
on the set of variables $x_1,\dots,x_m$, and then restrict the obtained
$S_m$-modules to the subgroup $\Gamma\subseteq S_m$.

%%%%%%%%%%%%%%%%%%%%%%%%%%%%%%%%%%%%%%%%
\subsection{Decomposition of $\mathcal A$ into irreducible 
$\Gamma$-modules}
%%%%%%%%%%%%%%%%%%%%%%%%%%%%%%%%%%%%%%%%

Recall that the action of $\Gamma$ on the edge coordinate functions is given
by a permutation of indices and is extended multiplicatively and linearly to
the whole algebra $\mathcal A$.
Since the action of $\Gamma$ preserves the degree of a monomial, each rank
component
$
A_k
$
is a $\Gamma$-invariant subspace, and therefore it is natural to pose the
problem of decomposing these spaces into irreducible components.

For each $0\le k\le m$, the vector space $A_k$ is the permutation
$S_m$-module on the set of all $k$-element subsets of $[m]$.
If $0\le k\le \lfloor m/2\rfloor$, then this module can be identified with
the standard permutation module $M^{(m-k,k)}$, realized as the vector space
with basis given by $(m-k,k)$-tabloids, see~\cite{Sagan}.
For $k>\lfloor m/2\rfloor$, the corresponding description is obtained from
the previous one by taking complements of subsets, that is, via the natural
isomorphism $A_k\cong A_{m-k}$.

Recall that a tabloid may be viewed as a Young tableau in which the
requirement that the entries be ordered within rows is removed; a tabloid is
determined only by the sets of entries in each row, independently of their
order. If
$0\le k\le \lfloor m/2\rfloor$, then $(m-k,k)$ is a partition of $m$, and
such tabloids canonically correspond to $k$-element subsets of $[m]$: the
second row of the tabloid gives the subset $S$, while the first row gives its
complement $[m]\setminus S$. Thus, the identification
$$
\boldsymbol{x}_S
\longleftrightarrow
\left\{\begin{array}{c}
[m]\setminus S\\
S
\end{array}\right\},
$$
defines an isomorphism of $S_m$-modules
$$
A_k\cong M^{(m-k,k)},
\qquad 0\le k\le \left\lfloor \frac m2\right\rfloor.
$$

Irreducible $S_m$-modules over a field of characteristic zero are parametrized
by partitions of the number~$m$. The module corresponding to a partition
$\lambda \vdash m$ is called a \emph{Specht module} and is denoted by
$S^\lambda$.
In general, the decomposition of the permutation module $M^\lambda$ into
irreducible submodules is described by \emph{Young's rule}
\cite[Theorem~2.11.2]{Sagan}, which in our case simplifies to
$$
M^{(m-k,k)}
\cong
\bigoplus_{j=0}^{\min(k,m-k)}
S^{(m-j,j)}.
$$

The dimension of the Specht module $S^\lambda$ is equal to the number
$f^\lambda$ of standard Young tableaux of shape $\lambda$ and is computed by
the hook-length formula, see for example \cite[Section~4.3]{Sagan}.
For the two-row partition $(m-k,k)$, the hook-length formula, see
Fulton~\cite[Exercise~9]{Fulton1997}, after simplification gives
\begin{equation} \label{dim_S}
\dim S^{(m-k,k)}=f^{(m-k,k)}
=
\frac{m!(m-2k+1)}{(m-k+1)!k!}
=
\binom{m}{k}-\binom{m}{k-1}.
\end{equation}
As we see, the dimension of the Specht module $S^{(m-k,k)}$ turns out to be
the same as the dimension of the space of primitive elements $P_k$. We shall
return to this fact in the next subsection.


\begin{thebibliography}{99}
\bibitem[Fulton1997]{Fulton1997} Fulton, W. (1997). Young tableaux: with applications to representation theory and geometry (No. 35). Cambridge University Press.
\bibitem[Sagan]{Sagan} Sagan, B. E. \emph{The Symmetric Group: Representations, Combinatorial Algorithms, and Symmetric Functions}. Graduate Texts in Mathematics, Vol. 203, Springer, 2001.
\end{thebibliography}
\end{document}
